\chapter*{Chapitre 5 — Distribution et émission de chaleur}
\addcontentsline{toc}{chapter}{Chapitre 5 — Distribution et émission de chaleur}

% ============================================================
\section*{Énoncés}
% ============================================================

\exercice{Réseau bitube — sélection des radiateurs d'un appartement}

Un appartement de 4 pièces doit être équipé de radiateurs à eau chaude 70/55\,°C ($T_{amb} = 20$\,°C, exposant $n = 1{,}3$). Les déperditions par pièce sont :

\begin{center}
\begin{tabular}{lcc}
  \toprule
  Pièce & Déperditions (W) & Surface (m²) \\
  \midrule
  Séjour & 1600 & 25 \\
  Chambre 1 & 800 & 12 \\
  Chambre 2 & 700 & 10 \\
  Cuisine & 500 & 8 \\
  \bottomrule
\end{tabular}
\end{center}

\begin{enumerate}
  \item Calculer la puissance nominale nécessaire pour chaque pièce.
  \item Calculer la puissance totale à fournir. Vérifier la cohérence avec le dimensionnement de la chaudière.
\end{enumerate}

% ----------------------------------------------------------

\exercice{Plancher chauffant — dimensionnement complet d'une pièce}

Un salon de 20\,m² présente des déperditions de \SI{900}{\watt}. Revêtement : parquet flottant compatible PCH ($R_{revêt} = \SI{0{,}10}{\square\metre\kelvin\per\watt}$). Le PCH fonctionne en 40/32\,°C, $T_{amb} = 20$\,°C.

\begin{enumerate}
  \item Calculer la puissance surfacique nécessaire.
  \item Calculer le coefficient $B$ et la puissance surfacique disponible avec $p = 150$\,mm.
  \item Calculer la longueur de tube nécessaire.
  \item Vérifier si une seule boucle suffit.
\end{enumerate}

% ----------------------------------------------------------

\exercice{Diagnostic d'un radiateur sous-performant}

Un technicien mesure sur un radiateur de bureau les valeurs suivantes :
\begin{itemize}
  \item Température entrée eau : 68\,°C
  \item Température sortie eau : 58\,°C
  \item Débit eau mesuré : 40\,L/h
  \item Température ambiante : 18\,°C
\end{itemize}
Le local présente des déperditions de 1200\,W. Le radiateur a une puissance nominale de 1800\,W (en 75/65/20\,°C, $n = 1{,}3$).

\begin{enumerate}
  \item Calculer la puissance thermique réellement délivrée.
  \item Calculer la puissance théorique attendue dans ces conditions.
  \item Analyser l'écart et proposer des pistes de diagnostic.
\end{enumerate}

\newpage
% ============================================================
\section*{Corrigés}
% ============================================================

\subsection*{Corrigé — Exercice 1}

\textit{Question 1 — Écart de température réel :}
\[
  T_{moy} = \frac{70 + 55}{2} = \SI{62{,}5}{\celsius}, \quad \Delta T_{réel} = 62{,}5 - 20 = \SI{42{,}5}{\kelvin}
\]
\[
  \text{Facteur correction} = \left(\frac{42{,}5}{50}\right)^{1{,}3} = (0{,}85)^{1{,}3} = 0{,}826
\]

\begin{methode}[title=Correction de puissance d'un émetteur]
La puissance nominale d'un radiateur est mesurée en conditions normalisées (75/65/20\,°C,
soit $\Delta T_{norm} = 50$\,K). Pour des conditions réelles différentes, on applique :
\[
  \Phi_{réel} = \Phi_{nom} \times \left(\frac{\Delta T_{réel}}{\Delta T_{norm}}\right)^n
\]
Pour dimensionner, on cherche la puissance nominale à sélectionner pour obtenir les déperditions en conditions réelles :
\[
  \Phi_{nom,néc} = \frac{\Phi_{dép}}{\text{facteur correction}}
\]
\end{methode}

\begin{center}
\begin{tabular}{lcc}
  \toprule
  Pièce & Déperditions (W) & $\Phi_{nom,néc}$ = Dép. / 0,826 (W) \\
  \midrule
  Séjour & 1600 & 1937 \\
  Chambre 1 & 800 & 969 \\
  Chambre 2 & 700 & 848 \\
  Cuisine & 500 & 605 \\
  \midrule
  \textbf{Total} & \textbf{3600} & \textbf{4359} \\
  \bottomrule
\end{tabular}
\end{center}

\textit{Question 2 :} La puissance totale de déperditions est 3600\,W. La chaudière devra fournir au minimum 3,6\,kW, majoré des pertes de distribution (5–10\,\%) → chaudière $\geq$ 4\,kW.

\[
  \boxed{P_{chaudière} \geq \SI{4}{\kilo\watt}}
\]

% ----------------------------------------------------------

\subsection*{Corrigé — Exercice 2}

\textit{Question 1 :}
\[
  q_{néc} = \frac{900}{20} = \SI{45}{\watt\per\square\metre}
\]

\[
  \boxed{q_{néc} = \SI{45}{\watt\per\square\metre}}
\]

\textit{Question 2 :}
\[
  B = \frac{1}{0{,}173 + 0{,}10} = \frac{1}{0{,}273} = \SI{3{,}66}{\watt\per\square\metre\per\kelvin}
\]
\[
  T_{moy} = \frac{40 + 32}{2} = \SI{36}{\celsius}, \quad \Delta T = 36 - 20 = \SI{16}{\kelvin}
\]
\[
  q_{disp} = B \times \frac{\Delta T}{p} = 3{,}66 \times \frac{16}{0{,}15} = \SI{390}{\watt\per\square\metre}
\]

La puissance disponible (390\,W/m²) est très largement supérieure à la puissance nécessaire (45\,W/m²). Même avec $p = 300$\,mm, $q = 195$\,W/m² >> 45\,W/m². Choisir $p = \SI{300}{\milli\metre}$.

\textit{Question 3 — Longueur de tube :}
\[
  L_{tube} = \frac{S}{p} = \frac{20}{0{,}30} = \SI{67}{\metre}
\]

\[
  \boxed{L_{tube} = \SI{67}{\metre}}
\]

\textit{Question 4 :} 67\,m < 100\,m → une seule boucle suffit.

% ----------------------------------------------------------

\subsection*{Corrigé — Exercice 3}

\textit{Question 1 — Puissance réelle mesurée :}
\[
  \dot{m} = \frac{40}{3600} \times 1000 \times \frac{1}{1000} = \SI{0{,}0111}{\kilogram\per\second}
\]
\[
  \dot{Q}_{réel} = \dot{m} \times c_p \times \Delta T = 0{,}0111 \times 4186 \times (68 - 58) = \SI{465}{\watt}
\]

\[
  \boxed{\dot{Q}_{réel} = \SI{465}{\watt}}
\]

\textit{Question 2 — Puissance théorique attendue :}
\[
  T_{moy} = \frac{68 + 58}{2} = \SI{63}{\celsius}, \quad \Delta T_{réel} = 63 - 18 = \SI{45}{\kelvin}
\]
\[
  \dot{Q}_{théo} = 1800 \times \left(\frac{45}{50}\right)^{1{,}3} = 1800 \times (0{,}90)^{1{,}3} = 1800 \times 0{,}870 = \SI{1566}{\watt}
\]

\[
  \boxed{\dot{Q}_{théo} = \SI{1566}{\watt}}
\]

\textit{Question 3 — Analyse :} Écart massif : 465\,W mesuré vs 1566\,W théorique (−70\,\%). Pistes de diagnostic :
\begin{itemize}
  \item Débit très faible (40\,L/h vs débit nominal probablement 100–150\,L/h) → robinet thermostatique bloqué ou partiellement fermé
  \item Présence d'air dans le radiateur → purge nécessaire
  \item Encrassement intérieur (boues) réduisant le débit et l'échange thermique
  \item Vanne d'équilibrage trop fermée → mesurer la pression différentielle
\end{itemize}
